\section{Análisis de Resultados}

Al cruzar los datos que medimos en el laboratorio con la teoría y los datos de catálogo del modelo LC1D0901 (Contactor 1), pudimos comprobar de primera mano cómo los aspectos constructivos del equipo afectan directamente su comportamiento eléctrico.

\subsection{Consumo en Mantenimiento y Eficiencia (Contactor 1)}
Según la hoja de instrucciones de Schneider Electric, los contactores pequeños de la línea TeSys D (del LC1D09 al LC1D15) cumplen con el Grado 3 de la norma de eficiencia GB 21518, lo que significa que en estado de mantenimiento (con el núcleo ya cerrado) no pueden consumir más de \qty{9,5}{\volt\ampere}.

Con los datos que medimos para el Contactor 1 (\qty{188}{\volt} y \qty{0,011}{\ampere}), calculamos su potencia aparente en régimen estable:
\[ S_{\text{mto, C1}} = \qty{188}{\volt} \times \qty{0,011}{\ampere} = \qty{2,07}{\volt\ampere} \]
Este valor (\qty{2,07}{\volt\ampere}) cumple perfectamente con la norma al estar muy por debajo del límite de \qty{9,5}{\volt\ampere}, y se acerca bastante a los \qty{7}{\volt\ampere} típicos que el fabricante declara en sus manuales comerciales para bobinas de corriente alterna. Esto nos demuestra que el circuito magnético de este contactor cerró de forma estanca y sin problemas.

\subsection{El Efecto del Entrehierro en la Corriente de Llamada y Mantenimiento}
Una de las observaciones más claras de la práctica fue la enorme diferencia entre la corriente de llamada ($I_{\text{llamada}}$) y la de mantenimiento ($I_{\text{mto}}$). Esto ocurre por la variación en la inductancia de la bobina debido al cambio en el entrehierro variable:
\begin{itemize}
    \item \textbf{Contactor 1 (LC1D0901):} La corriente de llamada (\qty{0,18}{\ampere}) resultó ser $\num{16,4}$ veces mayor que la de mantenimiento (\qty{0,011}{\ampere}). Al inicio, con el contactor abierto, el entrehierro es máximo, la inductancia es muy baja y la corriente se dispara. Cuando el martillo se pega al núcleo, el entrehierro se reduce a cero, la inductancia sube al máximo y la corriente cae al valor de régimen estable.
    \item \textbf{Contactor 2 (Modelo Grande):} Al ser un dispositivo de dimensiones mucho mayores (con un circuito magnético y contactos más pesados), demanda muchísima más energía para moverse. Por eso registramos una corriente de llamada de \qty{2,94}{\ampere} (un pico de potencia mucho más alto) y una corriente de mantenimiento de \qty{0,26}{\ampere}. A pesar de la escala, la proporción se mantiene: la corriente de llamada fue $\num{11,3}$ veces mayor que la de mantenimiento, validando el mismo principio físico de la variación de la inductancia por el cierre del entrehierro.
\end{itemize}

\subsection{Importancia del Entrehierro Permanente y las Espiras de Sombra}
Durante el desarmado del equipo, identificamos visualmente las espiras de sombra y los pequeños orificios en las caras del núcleo que forman el entrehierro permanente. Las mediciones eléctricas justifican por qué existen estas piezas:
\begin{itemize}
    \item Las espiras de sombra evitan que el contactor vibre o zumbe debido al paso por cero de la corriente alterna, manteniendo una fuerza de atracción constante. Durante las pruebas para el Contactor 1, se obtuvo que no se escuchaban vibraciones en estado estable, lo que indica que están cumpliendo su función. Por el contrario, en el Contactor 2 sí se registraron vibraciones audibles aun en estado estable; esto puede deberse a su tamaño mucho más grande o a un desgaste en su respectiva espira de sombra.
    \item El entrehierro permanente evita que el magnetismo remanente deje ``pegado'' el contactor cuando se corta la energía. En el Contactor 2, pudimos medir que al bajar la tensión a \qty{54}{\volt} (tensión mínima de operación/caída), el resorte vence la fuerza magnética residual y abre los contactos de golpe. Esto demuestra que el entrehierro permanente rompe la retención magnética residual con éxito para garantizar una apertura rápida y segura.
\end{itemize}

\subsection{Comparación de los Niveles de Tensión}
Al analizar las tensiones críticas, se nota claramente la diferencia de diseño entre ambos equipos. El Contactor 1 requirió una tensión de mantenimiento alta (\qty{188}{\volt}), lo que concuerda con una bobina diseñada para operar en rangos industriales estándar de \qty{220}{\volt}.

Por su parte, el Contactor 2 se mantuvo acoplado firmemente con \qty{96}{\volt} y abrió sus contactos al bajar de \qty{54}{\volt}. Estos niveles de tensión de control, sumados a su alto consumo en amperios, nos indican que este contactor grande utiliza una bobina diseñada para rangos de voltaje menores (o un circuito de mando intermedio), donde se prioriza una alta circulación de corriente para generar el enorme flujo magnético necesario para mover su estructura.

\subsection{Efecto de un Obstáculo en el Entrehierro}
Durante los ensayos prácticos, se introdujo un obstáculo (consistente en una hoja de papel doblada) en el entrehierro del contactor para evaluar su impacto en las variables eléctricas y mecánicas. Se observó que, ante la presencia del papel que impedía el cierre magnético estanco, la corriente de mantenimiento aumentó notablemente hasta alcanzar un valor de \qty{0,40}{\ampere}, y de forma simultánea se incrementó significativamente la vibración mecánica y el zumbido emitido por el dispositivo.

Este comportamiento tiene una clara explicación electromagnética y mecánica: al obstaculizar el cierre completo del martillo contra el núcleo, se introduce un entrehierro artificial que incrementa la reluctancia del circuito magnético. Al aumentar la reluctancia, la inductancia de la bobina disminuye considerablemente, lo que a su vez reduce la reactancia inductiva. Por consiguiente, bajo una tensión constante, la corriente de mantenimiento aumenta significativamente para intentar compensar la reluctancia y mantener el flujo magnético necesario para la retención del núcleo. Asimismo, el aumento drástico en las vibraciones ocurre porque la fuerza de atracción magnética disminuye con la distancia del entrehierro y las espiras de sombra pierden efectividad al no existir un acoplamiento directo y hermético de las caras ferromagnéticas, permitiendo que la fuerza del resorte y la oscilación de la corriente alterna generen un golpeteo constante y ruidoso contra el obstáculo.
